\textbf{(i)} A function on isomorphism classes extends uniquely to the free abelian group on those classes. Additivity says exactly that the defining relations lie in its kernel. It therefore factors uniquely through $K(A)$.

\textbf{(ii)} By \exref{7.18}, a finitely generated module has a finite filtration with factors $A/\mathfrak{p}_i$. The short exact sequences of successive terms give $\gamma_A(M)=\sum_i\gamma_A(A/\mathfrak{p}_i)$.

\textbf{(iii)} If $\mathfrak{p}\ne0$ is prime in a principal ideal domain, write $\mathfrak{p}=(a)$. The exact sequence $0\to A\xrightarrow{a}A\to A/\mathfrak{p}\to0$ gives $\gamma_A(A/\mathfrak{p})=0$. Thus (ii) makes $K(A)$ cyclic, generated by $\gamma_A(A)$. The function $M\mapsto\dim_{\operatorname{Frac}(A)}(\operatorname{Frac}(A)\otimes_A M)$ is additive because localization is exact. It induces a homomorphism $K(A)\to\mathbb{Z}$ taking $\gamma_A(A)$ to $1$, so $K(A)\cong\mathbb{Z}$.

\textbf{(iv)} Since $B$ is finite over $A$, a finitely generated $B$-module is finitely generated over $A$. Restriction of scalars preserves exact sequences, so (i) gives $f_!$ with $f_!(\gamma_B(M))=\gamma_A(M)$. Restricting from $C$ to $A$ directly or through $B$ gives the same module, proving $(g\circ f)_!=f_!\circ g_!$.
